\setcounter{page}{1}
\setcounter{figure}{0}
\setcounter{chapter}{2}

%\setcounter{table}{0}
%\thispagestyle{empty}
%\pagestyle{empty}
%\fancyfoot{}
%\fancyfoot[LE,RO]{}
\chapter{Cantidades Cinemáticas}

\vspace{-15ex}
\begin{flushright}
\linea{2}
\end{flushright}
\raggedright
\resp{.75\textwidth}{Los n\'umeros entre par\'entesis hacen referencia al libro principal de consulta. Por ejemplo, \textcolor{c5}{[2.1]} es el ejercicio \#1 del Cap\'itulo 2.}{white}{noticias}
\begin{enumerate}
\item \textcolor{c5}{[2.1]} Un automóvil viaja en la dirección $+x$ sobre un camino recto y
nivelado. En los primeros 4.00 s de su movimiento, la velocidad
media del automóvil es $v_{\text{med}-x} = 6.25 \text{ m/s}$. ¿Qué distancia viaja el automóvil en 4.00 s?

\item \textcolor{c5}{[2.4]} \textbf{De pilar a poste}. Partiendo de un pilar, usted corre 200 m al este (en la dirección $+x$) con rapidez media de 5.0 m/s, luego
280 m al oeste con rapidez media de 4.0 m/s hasta un poste. Calcule a) su rapidez media del pilar al poste y b) su velocidad media del pilar al poste.

\item \textcolor{c5}{[2.6]} Un Honda Civic viaja en línea recta en carretera. Su distancia $x$ a partir de un letrero de alto está dada en función del tiempo $t$ por
la ecuación $x(t) = \alpha t^
2 - \beta t^
3$, donde $\alpha = 1.50 \text{ m/s}^2$ y $\beta = 0.0500 \text{ m/s}^3$.
Calcule la velocidad media del automóvil para los intervalos a) $t = 0$
a $t = 2.00 \text{s}$; b) $t = 0$ a $t = 4.00 \text{ s}$; c) $t = 2.00 \text{ s}$ a $t = 4.00 \text{ s}$.

\item \textcolor{c5}{[2.7]} Un automóvil está detenido ante un semáforo. Después,
viaja en línea recta y su distancia con respecto al semáforo está dada
por $x(t) = bt^2 - ct^3$, donde $b = 2.40 \text{ m/s}^2$ y $c = 0.120 \text{ m/s}^3$. a) Calcule la velocidad media del automóvil entre el intervalo $t = 0$ a $t = 10.0 \text{ s}$.
b) Calcule la velocidad instantánea del automóvil en $t = 0$, $t = 5.0 \text{ s}$ y
$t = 10.0 \text{ s}$. c) ¿Cuánto tiempo después de que el auto arrancó vuelve
a estar detenido?


\item \textcolor{c5}{[2.18]} La posición del parachoques (defensa) frontal de un automóvil de pruebas controlado por un microprocesador está dada por $x(t) = 2.17 \text{ m} + (4.80 \text{ m/s}^2)t^2 - (0.100 \text{ m/s}^6)t^6$. a) Obtenga su
posición y aceleración en los instantes en que tiene velocidad cero. b) Dibuje las gráficas $x\text{ vs }t$, $v_x\text{ vs }t$ y $a_x\text{ vs }t$ para el movimiento del frente
del auto entre $t = 0$ y $t = 2.00 \text{ s}$. 
\end{enumerate}


